\section{Thứ tự của hai phép, đo ở lúc khởi tạo}
\label{sec:ch08-thu-tu-ln}

\index{layer normalization!post-LN và pre-LN}%
Mục~\ref{sec:ch07-quanh-tang} đặt \eqref{eq:ch07-post-ln} xuống và hẹn chương
này. Bài chuẩn hóa \emph{sau} phép cộng:

\begin{equation}
  \vect{x}_{l+1} = \layernorm(\vect{x}_l + \mathrm{Sublayer}(\vect{x}_l)),
  \qquad
  \vect{x}_{l+1} = \vect{x}_l + \mathrm{Sublayer}(\layernorm(\vect{x}_l)).
  \label{eq:ch08-hai-thu-tu}
\end{equation}

Vế trái là post-LN, cách bài viết. Vế phải là pre-LN, và hai chữ ấy là tên gọi
về sau chứ không phải chữ của bài. Cùng ba phép, cùng số tầng, khác chỗ đặt một
dấu ngoặc.

\index{kết nối tắt!đường tắt đi qua hay không đi qua chuẩn hóa}%
Chỗ khác nhau nằm ở đường tắt. Trong vế phải, \(\vect{x}_l\) đi thẳng sang
\(\vect{x}_{l+1}\) bằng một phép cộng và không gặp phép chuẩn hóa nào; xếp
\(L\) tầng như thế thì có một đường từ đầu vào tới đầu ra chỉ gồm các phép cộng.
Trong vế trái thì không: mỗi tầng, đường tắt chui qua một layer norm.

Đó đúng là hình dạng mục~\ref{sec:ch07-quanh-tang} đã dẫn ra cho hộp cầu nối He
2015, và đúng hình dạng của vòng quay lỗi ở chương~\ref{ch:lstm}. Hai chương
trước đã dựng sẵn công cụ để đọc nó. Hình~\ref{fig:ch08-thu-tu-ln} vẽ cả hai.

\begin{figure}[htbp]
  \centering
  % Thứ tự cộng và LayerNorm trong một tầng con: post-LN so với pre-LN.
% Vẽ đúng theo Xiong và cộng sự 2020, bảng 1:
%   post-LN (Vaswani và cộng sự dùng): x_{l+1} = LayerNorm(x_l + Sublayer(x_l))
%   pre-LN:                             x_{l+1} = x_l + Sublayer(LayerNorm(x_l))
% Trọng tâm của hình là đường tắt residual từ x_l tới x_{l+1}: vẽ đậm ở cả hai
% khối để mắt bắt vào nó trước, rồi để người đọc tự thấy nó đi qua LayerNorm ở
% post-LN (trái) nhưng không bao giờ chạm LayerNorm ở pre-LN (phải). Pre-LN
% còn cần một LayerNorm thêm, nằm ngoài khối lặp N tầng, sau tầng cuối cùng
% x_L - vẽ tách hẳn phía trên khung chấm chấm và tô nền xám nhạt để đánh dấu
% nó không phải một bước lặp lại như phần còn lại.
\begin{tikzpicture}[
  >= Stealth,
  hop/.style      = {draw, rectangle, inner sep = 2pt, font = \scriptsize,
                      align = center},
  cong/.style     = {draw, circle, minimum size = 4mm, inner sep = 0pt,
                      font = \scriptsize},
  ngoai/.style    = {fill = black!10},
  tinh/.style     = {->, line width = 0.6pt, black!75},
  duongtat/.style = {->, line width = 1.6pt, black},
  khung/.style    = {black!45, densely dotted, line width = 0.5pt},
  nho/.style      = {font = \scriptsize},
  ghichu/.style   = {font = \footnotesize, black!60},
]

  % ================= POST-LN (trái, x = 0) =================
  \node[nho, anchor = east] at (-0.15, 0) {\(\vect{x}_l\)};
  \draw[tinh] (0, 0.15) -- (0, 0.6);
  \node[hop, minimum width = 2.6cm, minimum height = 0.8cm]
    (sub_a) at (0, 1.0) {Sublayer};
  \draw[tinh] (0, 1.4) -- (0, 1.85);
  \node[cong] (add_a) at (0, 2.05) {\(+\)};
  \draw[duongtat] (0, 2.25) -- (0, 2.7);
  \node[hop, minimum width = 2.6cm, minimum height = 0.8cm]
    (ln_a) at (0, 3.1) {LayerNorm};
  \draw[duongtat] (0, 3.5) -- (0, 3.95);
  \node[nho, anchor = south] at (0, 3.95) {\(\vect{x}_{l+1}\)};

  % đường tắt: từ x_l, vòng qua bên trái hộp Sublayer, vào cạnh tây của phép
  % cộng, rồi đi tiếp - vẫn đậm - xuyên qua LayerNorm lên tới x_{l+1}. Đoạn
  % "đi tiếp" đã được vẽ ở trên (hai mũi tên duongtat quanh add_a và ln_a);
  % đoạn vòng bên trái là phần còn thiếu.
  \draw[duongtat] (0, 0.15) -- (-1.5, 0.15) -- (-1.5, 2.05) -- (add_a.west);

  \begin{scope}[on background layer]
    \draw[khung, rounded corners = 2pt] (-1.9, -0.4) rectangle (1.9, 4.3);
  \end{scope}
  \node[ghichu, anchor = west] at (2.05, 1.95) {N x};

  \node[nho, align = center, text width = 3.0cm] at (0, 5.5)
    {không cần LayerNorm thêm: mỗi tầng đã kết thúc bằng LayerNorm};

  \node[font = \bfseries\scriptsize] at (0, -0.85) {post-LN};
  \node[nho, black!70] at (0, -1.35)
    {\(\vect{x}_{l+1} = \layernorm(\vect{x}_l + \mathrm{Sublayer}(\vect{x}_l))\)};

  % ================= PRE-LN (phải, x = 6.2) =================
  \node[nho, anchor = east] at (6.05, 0) {\(\vect{x}_l\)};
  \draw[tinh] (6.2, 0.15) -- (6.2, 0.6);
  \node[hop, minimum width = 2.6cm, minimum height = 0.8cm]
    (ln_b) at (6.2, 1.0) {LayerNorm};
  \draw[tinh] (6.2, 1.4) -- (6.2, 1.85);
  \node[hop, minimum width = 2.6cm, minimum height = 0.8cm]
    (sub_b) at (6.2, 2.25) {Sublayer};
  \draw[tinh] (6.2, 2.65) -- (6.2, 3.1);
  \node[cong] (add_b) at (6.2, 3.3) {\(+\)};
  \draw[duongtat] (6.2, 3.5) -- (6.2, 3.95);
  \node[nho, anchor = south] at (6.2, 3.95) {\(\vect{x}_{l+1}\)};

  % đường tắt: từ x_l, vòng qua bên phải cả LayerNorm lẫn Sublayer, vào cạnh
  % đông của phép cộng - không chạm hộp nào, đúng chỗ khác với post-LN
  \draw[duongtat] (6.2, 0.15) -- (7.7, 0.15) -- (7.7, 3.3) -- (add_b.east);

  \begin{scope}[on background layer]
    \draw[khung, rounded corners = 2pt] (4.3, -0.4) rectangle (8.1, 4.3);
  \end{scope}
  \node[ghichu, anchor = west] at (8.25, 1.95) {N x};

  % --- LayerNorm thêm, ngoài khối lặp: một lần duy nhất sau tầng cuối x_L,
  % trước khi sang bộ dự đoán. Đây là chỗ cài đặt dễ quên nhất của pre-LN.
  \draw[tinh] (6.2, 4.3) -- (6.2, 4.75);
  \node[nho, anchor = west] at (6.35, 4.75) {\(\vect{x}_L\)};
  \draw[tinh] (6.2, 4.75) -- (6.2, 5.15);
  \node[hop, ngoai, minimum width = 2.6cm, minimum height = 0.8cm]
    (lnfinal) at (6.2, 5.55) {LayerNorm};
  \draw[tinh] (6.2, 5.95) -- (6.2, 6.4);
  \node[nho, anchor = south] at (6.2, 6.4) {tới Linear và softmax};

  \node[ghichu, anchor = west, align = left, text width = 3.2cm]
    at (8.25, 5.55)
    {ngoài khối N tầng: một lần, sau tầng thứ L, thêm \(4\dmodel\) tham số};

  \node[font = \bfseries\scriptsize] at (6.2, -0.85) {pre-LN};
  \node[nho, black!70] at (6.2, -1.35)
    {\(\vect{x}_{l+1} = \vect{x}_l + \mathrm{Sublayer}(\layernorm(\vect{x}_l))\)};

  % ================= chú giải đường đậm, đặt dưới cả hai khối =================
  \node[nho, anchor = north, align = left, black!70, text width = 11cm]
    at (3.1, -2.0)
    {đường đậm: đường tắt residual từ \(\vect{x}_l\) tới \(\vect{x}_{l+1}\). Ở
     post-LN nó buộc phải đi qua LayerNorm; ở pre-LN thì không bao giờ chạm
     LayerNorm.};

\end{tikzpicture}

  \caption{Hai thứ tự của \eqref{eq:ch08-hai-thu-tu}. Đường đậm là đường tắt:
    bên post-LN nó chui qua ô chuẩn hóa ở mỗi tầng, bên pre-LN nó đi thẳng từ
    đầu vào tới đầu ra của cả chồng tầng. Ô rời phía dưới bên phải là cái layer
    norm cuối cùng mà pre-LN bắt buộc phải có và rất dễ quên, vì nó nằm ngoài
    khối lặp lại chứ không nằm trong.}
  \label{fig:ch08-thu-tu-ln}
\end{figure}

\begin{bridge}{Xiong và cộng sự 2020: vì sao mục 5.3 của bài phải tồn tại}
\index{layer normalization!Xiong 2020}%
Bài này chứng minh rằng thứ tự ở \eqref{eq:ch08-hai-thu-tu} là lý do lịch warmup
của mục 5.3 phải có~\autocite{xiong2020layernorm}. Định lý 1 chặn chuẩn
Frobenius của gradient trên ma trận thứ hai của mạng truyền thẳng ở tầng cuối:

\begin{quote}
\enquote{for the Post-LN Transformer with \(L\) layers, the gradient of the
parameters of the last layer satisfies
\(\norm{\partial \tilde{\mathcal{L}} / \partial W^{2,L}}_F \leq
\mathcal{O}(d\sqrt{\ln d})\) and for the Pre-LN Transformer with \(L\) layers,
\(\norm{\partial \tilde{\mathcal{L}} / \partial W^{2,L}}_F \leq
\mathcal{O}(d\sqrt{\ln d / L})\).}
\end{quote}

Upper bound của post-LN không chứa \(L\); của pre-LN thì nhỏ đi theo
\(1/\sqrt{L}\).
Từ đó bài rút ra kết luận nêu ngay trong tóm tắt:
\enquote{This motivates us to remove the warm-up stage for the training of
Pre-LN Transformers.}

\textbf{Định lý ấy chứng minh trong một cấu hình rút gọn, và ba phép rút gọn ấy
đáng đọc trước khi trích nó.} Mục 3.3 liệt kê: một head thay vì nhiều head; khởi
tạo \(\mat{W}^Q\) và \(\mat{W}^K\) bằng \emph{ma trận không}, nên ở lúc khởi tạo
attention là một trung bình cộng đều; và đầu vào lấy từ cùng một phân phối
Gauss. Phép thứ hai là phép nặng nhất: định lý nói về một mô hình mà cơ chế
attention đã tắt. Khởi tạo dùng trong chứng minh là Xavier, không phải mặc định
của \code{nn.Linear} mà mục~\ref{sec:ch07-chia-can} đo được là phương sai 1/3.

Số đo của chính bài, trên IWSLT14 Đức-Anh với post-LN và Adam: bỏ warmup thì
Bilingual Evaluation Understudy, viết tắt là BLEU, chỉ đạt 8.45, còn có warmup
thì \enquote{around 34}. Đặt
\(T_{\mathrm{warmup}} = 500\) cho 31.16 và 2.77 ở hai learning rate khác nhau.
Bốn con số ấy là của bài, không phải của tôi.

Một chi tiết thư mục: bản arXiv của bài ghi sai venue của chính nó. Chân trang
1 in \enquote{PMLR 108, 2020} cho ICML lần thứ 37, mà ICML 37 là PMLR tập 119.
Mục lục của PMLR mới là nguồn đúng, và \code{refs.bib} lấy số tập từ đó.
\end{bridge}

\subsection*{Cái layer norm cuối cùng, và vì sao nó dễ mất}

\index{layer normalization!pre-LN cần một chuẩn hóa cuối}%
Bảng 1 của Xiong dành hẳn một dòng riêng cho một chi tiết mà tôi suýt bỏ qua khi
cài: pre-LN cần thêm một layer norm sau cả chồng tầng, trước phép dự đoán.

Lý do thì hiển nhiên khi đã viết \eqref{eq:ch08-hai-thu-tu} ra. Mỗi tầng pre-LN
chuẩn hóa \emph{đầu vào} của chính nó, nên đầu ra của tầng cuối chưa được chuẩn
hóa lần nào. Bỏ dòng ấy đi thì mô hình vẫn chạy, vẫn học, mất mát vẫn giảm, và
cái đi vào softmax là một vector không ai chuẩn hóa. Đây đúng họ với những lỗi
mục~\ref{sec:ch07-che} đã gặp: không exception, không cảnh báo.

Nó cũng là thứ duy nhất phân biệt được hai thứ tự từ bên ngoài.

\begin{measured}{số tham số của hai thứ tự, cùng một cấu hình}
\begin{minted}{text}
d_model  post-LN params  pre-LN params  difference  4*d_model
24       35,845          35,941         96          96
64       238,885         239,141        256         256
512      14,755,877      14,757,925     2048        2048
\end{minted}

2 tầng mỗi phía. Đo bằng \code{experiments/ch08\_norm.py} tại tag \code{ch08}.
\end{measured}

Đúng \(4\dmodel\): một weight và một bias, trên mỗi trong hai chồng tầng. Cột
post-LN ở hàng đầu là 35845, đúng con số mục~\ref{sec:ch07-corpus} in ra.

\subsection*{Đo định lý trên mô hình thật}

\index{layer normalization!gradient lúc khởi tạo theo độ sâu}%
Định lý là một phát biểu về lúc khởi tạo, nên kiểm nó không cần huấn luyện gì
cả: dựng mô hình, đẩy một batch qua, gọi \code{backward}, đọc chuẩn của gradient.
Rẻ, và tất định.

Một chỗ phải cẩn thận về \emph{tầng cuối} nghĩa là tầng nào. Định lý nói về tầng
mà softmax đọc đầu ra. Trong một kiến trúc encoder-decoder thì đó là tầng cuối
của \emph{decoder}, không phải của encoder: gradient tới tầng cuối encoder phải
đi ngược qua toàn bộ chồng decoder trước, nên nó không phải đại lượng định lý
chặn. Bảng dưới in cả hai để thấy chúng khác nhau.

\begin{measured}{chuẩn Frobenius của gradient ở lúc khởi tạo, theo độ sâu}
\begin{minted}[fontsize=\footnotesize]{text}
L    post-LN       pre-LN        post/pre   post enc last  pre enc last
1    0.349135      0.077308      4.52       0.103942       0.004518
2    0.268458      0.056034      4.79       0.075832       0.003047
4    0.272682      0.059755      4.56       0.140697       0.006837
8    0.291909      0.057998      5.03       0.160825       0.007671
16   0.353580      0.060844      5.81       0.530925       0.011664
\end{minted}

\(\dmodel = 24\), 4 head, một batch 128, không huấn luyện.

Đo bằng \code{experiments/ch08\_norm.py} tại tag \code{ch08}.
\end{measured}

Chia mỗi cột cho chính hàng \(L = 1\) của nó thì đọc được hình dạng:

\begin{measured}{cùng số liệu, chuẩn hóa theo hàng \(L = 1\)}
\begin{minted}{text}
L    post/post(L=1)  pre/pre(L=1)  1/sqrt(L)
1    1.0000          1.0000        1.0000
2    0.7689          0.7248        0.7071
4    0.7810          0.7729        0.5000
8    0.8361          0.7502        0.3536
16   1.0127          0.7870        0.2500
\end{minted}

Đo bằng \code{experiments/ch08\_norm.py} tại tag \code{ch08}.
\end{measured}

\index{layer normalization!nửa nào của định lý tái lập được}%
Bảng này nói ba chuyện và tôi tách chúng ra, vì gộp lại là đọc quá tay.

\textbf{Nửa thứ nhất của định lý tái lập được.} Cột post-LN không đổi theo độ
sâu: 0.349135 ở \(L = 1\) và 0.353580 ở \(L = 16\), tỷ lệ 1.0127 sau khi đi qua
mười sáu tầng. Định lý nói upper bound của post-LN không chứa \(L\), và số đo đứng yên
đúng như thế.

\textbf{Nửa thứ hai thì không.} Cột pre-LN đáng lẽ theo \(1/\sqrt{L}\), tức còn
0.2500 ở \(L = 16\). Đo được 0.7870. Nó cũng gần như đứng yên.

Tôi không kết luận định lý sai, và có ba lý do khác nhau mà tôi không tách được
bằng số đo này. Định lý là một \emph{upper bound}, không phải một đẳng thức, nên
một phép đo nằm yên dưới một cái trần đang tụt không mâu thuẫn với nó, chỉ nói
rằng cái trần chưa chạm tới. Cấu hình rút gọn ở hộp cầu nối không phải mô hình
này, và phép đặt \(\mat{W}^Q = \mat{W}^K = 0\) đổi hẳn thứ đang được lấy đạo
hàm. Và \(\dmodel = 24\) với \(L\) tới 16 là một mô hình rất hẹp và rất sâu, tức
xa cấu hình mà cả bài lẫn thực tế quan tâm.

\textbf{Cái đứng vững nhất là tỷ số.} Cột \code{post/pre} nằm giữa 4.52 và 5.81
ở mọi độ sâu: gradient của post-LN ở đúng chỗ ấy lớn hơn của pre-LN khoảng năm
lần. Đó là nội dung thực dụng của định lý, và nó không phụ thuộc vào việc nửa
thứ hai có tái lập hay không. Một gradient lớn gấp năm ở ngay lúc bắt đầu là thứ
mà một learning rate cố định phải sống chung, và mục~\ref{sec:ch08-warmup} là
chỗ nói cách bài sống chung với nó.

Hai cột encoder thì kể chuyện ngược lại và đáng một câu: chúng \emph{tăng} theo
độ sâu, post-LN từ 0.103942 lên 0.530925, gấp 5.1 lần. Không mâu thuẫn với gì
cả, vì đó không phải đại lượng định lý chặn; nó là gradient đã đi ngược qua cả
chồng decoder rồi mới tới nơi. Tôi in ra vì một bảng chỉ có cột decoder sẽ để
người đọc tưởng chuyện độ sâu là chuyện đã xong.
